\paragraph{\trEquivariant\trFukayaCategory{}。}\label{sub:equivariant-fukaya-category}
现在，把我们讨论的两条线索合并起来。设 $M$ 是一个\trCorneredBoundary{}的\trExactSymplectic\trManifold，其上装备着一个服从\cref{thm:fixed-point-set-assumption}的\trInvolution{}$\iota_M$。此外，它还应带有一个不变的二次复\trVolume{}形式 $\eta_M^2$。
考虑 $M$ 上的预\trEquivariant\trExact\trLagrange\trBrane{}，并附加如下条件：若它们是\trEquivariant{}的，则\cref{thm:normal-angle-assumption}必须被满足（除非它们与\trFixedPoint{}集不相交，这就蕴含一个\trNormalAngle{}的选取）。为任何一对这样的\trBrane{}选好\trFloerDatum{}。这应当相对于 $\ZZ/2$-作用 $\epsilon$ 是自然的，即
\begin{equation}\label{eq:pulledback-floer-datum}
(H_{L^{\#}_{0,-},L^{\#}_{1,-}},J_{L^{\#}_{0,-},L^{\#}_{1,-}})
=(\iota_M^*H_{L^{\#}_{0,+},L^{\#}_{1,+}},\iota_M^*J_{L^{\#}_{0,+},L^{\#}_{1,+}}).
\end{equation}
如果所涉及的两个\trBrane{}都是\trEquivariant\trBrane{}，这就意味着\trFloerDatum{}必须是 $\iota_M$-不变的；在这种情形下，我们再按\cref{thm:floer-datum-assumption}进一步约束选取。我们还要求所有\trFloerDatum{}都是\trRegular{}的。对\trEquivariant\trBrane{}的对，我们在\cref{thm:regularity-fixed-set}之后看到这可以做到；而在所有其他情形都不成问题，因为 \eqref{eq:pulledback-floer-datum} 把一个\trFloerDatum{}关联到另一个：选两者之一为\trRegular{}的，另一个显然也会具有同样的性质。\trFloerDatum{}选取中的这种对称性意味着，我们有\trCanonical{}的\trIsomorphism
\begin{equation}\label{eq:epsilon-chain-isomorphism}
\epsilon\colon CF^*(L^{\#}_{0,+},L^{\#}_{1,+})\longrightarrow CF^*(L^{\#}_{0,-},L^{\#}_{1,-}),
\end{equation}
它们与\trDifferential{}严格交换，而且其平方为恒等。接下来，为\trObject{}的全部 $(d{+}1)$-元组选好\trPerturbationDatum{}，既满足通常的\trConsistency{}要求，也满足 \eqref{eq:pulledback-floer-datum} 的类似物。多亏\cref{thm:regularity-polygons}，在这一类之内可以达到\trRegularity{}，而所得的诸\trComposition{}映射 $\mu^d$ 与 \eqref{eq:epsilon-chain-isomorphism} 严格交换。其结果是，我们得到一个\trACategory{}$\scr{F}(M)^{\mathrm{naive}}$，即\textbf{\trEquivariant\trFukayaCategory}（equivariant Fukaya category），其\trObject{}是受上述条件约束的预\trEquivariant\trBrane{}，\trMorphism{}空间为 $CF^*(L^{\#}_{0,+},L^{\#}_{1,+})$。这个\trACategory{}支持一个由 \eqref{eq:epsilon-chain-isomorphism} 给出的朴素 $\ZZ/2$-作用。为把它与之前的构造相比较，让我们过渡到\trCohomology，那里显然有一个\trFullyFaithful{}的\trEmbedding{}$HF(M)^{\mathrm{naive}}\to HF(M)^{\mathrm{strict}}$，它与严格的 $\ZZ/2$-作用\trCompatible{}。这实际上是一个\trEquivalence{}：违反\cref{thm:normal-angle-assumption}中那些条件的不变的\trLagrange\trSubmanifold{}不能实现为 $HF(M)^{\mathrm{naive}}$ 中的\trEquivariant\trBrane{}，但可以把它们接纳为预\trEquivariant\trBrane{}——在\trCohomology{}层面上它们是\trIsomorphic{}的\trObject{}。换言之，这样的\trSubmanifold{}作为 $HF(M)^{\mathrm{naive}}$ 的\trObject{}是存在的，但无法把它们做成这个 $\ZZ/2$-作用的\trFixedPoint{}。由此得到的好处是：$HF(M)^{\mathrm{naive}}$ 的确是一个带朴素 $\ZZ/2$-作用的\trACategory{}的\trCohomology{}\trCategory{}（不像 $HF(M)^{\mathrm{strict}}$——那是一个纯粹\trCohomology{}层面的构造）。
为把 \eqref{eq:epsilon-chain-isomorphism} 描述得稍微详细一点，取某个 $y_+\in\scr{C}(L^{\#}_{0,+},L^{\#}_{1,+})$（我们在记号中把\trBraneStructure{}也标出来，因为\trHamiltonian{}量 $H$ 从而集合 $\scr{C}$ 现在依赖于它们），并考虑与之相伴的那对\trLinear\trLagrange\trBrane{}$\Lambda^{\#}_{k,+}$，即 $L^{\#}_{k,+}$ 在 $y_+(k)$ 处的\trTangentSpace{}。对于 $\scr{C}(L^{\#}_{0,-},L^{\#}_{1,-})$ 中相应的元素——即 $y_-=\iota_M(y_+)$——按\trObject{}上 $\ZZ/2$-作用的定义，与之相关的\trLinear\trLagrange\trBrane{}$\Lambda^{\#}_{k,-}$ 就是 $L^{\#}_{k,-}$ 的\trTangentSpace{}。作为我们\trEquivariant\trBraneStructure{}的一部分，我们有\trIsomorphism
\begin{equation}\label{eq:linear-brane-tensor-iso}
(f_{k,-})_{y_-(k)}\colon\Lambda^{\#}_{k,-}\longrightarrow
\Lambda^{\#}_{k,+}\otimes\lambda^{\mathrm{top}}(TL_{k,+})_{y_+(k)}
\end{equation}
\begingroup\tolerance=9999\emergencystretch=3em\hbadness=10000
\trCovering{}$D\iota_M\colon(TL_{k,-})_{y_-(k)}\longrightarrow(TL_{k,+})_{y_+(k)}$ 的作用。鉴于\cref{thm:grading-relative-orientation}，我们有优先的\trIsomorphism{}$\lambda^{\mathrm{top}}(TL_{k,+})_{y_+(k)}\cong\xi_{y_+(k)}$，并且还可以进一步利用沿 $y$ 的\trParallelTransport{}来等同 $\xi_{y_+(0)}\cong\xi_{y_+(1)}$。把所涉及的两个\trLinear\trLagrange\trBrane{}都与同一条实直线作\trTensorProduct{}，既不改变绝对\trIndex{}，也不改变（在相差\trCanonical\trIsomorphism{}的意义下）\trOrientation{}空间。于是，\eqref{eq:linear-brane-tensor-iso} 诱导一个\trIsomorphism{}$o(y_-)\to o(y_+)$。这些\trIsomorphism{}之逆的 $\KK$-\trNormalization{}构成 \eqref{eq:epsilon-chain-isomorphism}。这确是一个 $\ZZ/2$-作用这一事实，归结为 \eqref{eq:equivariant-brane-square} 中的条件以及如下事实：若用非平凡元素 $-e\in\ker(Pin_n\to O_n)$ 同时作用在两个\trLinear\trLagrange\trBrane{}$\Lambda^{\#}_{k,+}$ 上，则 $o(y)$ 上诱导的\trAutomorphism{}是恒等。\par\endgroup
有一种情形或许值得特别留意：设 $L^{\#}_k$ 是两个\trEquivariant\trBrane{}，而 $y$ 是 $\scr{C}(L^{\#}_0,L^{\#}_1)$ 中在 $\iota_M$ 作用下不变的元素。此时，上面的构造给出 $o(y)$ 的一个\trAutomorphism{}，其符号由局部几何数据确定。我们如同\cref{thm:equivariant-index-orientation}的证明那样作 $H=0$ 这一简化假设，于是 $y$ 是一个不变交点。回想由\cref{thm:codim-one-iso-coorientations}，诸\trIsomorphism{}$f_k$ 确定了不动部分 $L_k^{\mathrm{fix}}$ 的\trNormalBundle{}的优先\trTrivialization{}。在 $D\pi$ 之下，这些\trTrivialization{}都\trProjection{}为直线 $e^{j\sigma_k}\RR$ 的\trOrientation。若所得\trOrientation{}与由\trGenerator{}$e^{j\sigma_k}$ 给出的那个一致，就令 $w_k=+1$；否则令 $w_k=-1$。那么，
\begin{lemma}\label{thm:epsilon-invariant-part}
$|o(y)|_\KK$ 属于 $CF^*(L_0^{\#},L_1^{\#})$ 的 $\epsilon$-不变部分，若 $w_0w_1=+1$；否则属于反不变部分。
\end{lemma}
\begin{proof}
我们依循\cref{thm:index-free-signs-rephrase}的策略，选一条路径 $\lambda\in\scr{P}^-\,Gr(V)$，从 $\lambda(0)=(TL_0)_y$ 到 $\lambda(1)=(TL_1)_y$，它完全由 $\ZZ/2$-不变的\trLagrange{}子空间组成。以 $F_\lambda$ 记 $[0;1]$ 上相伴的\trVectorBundle{}。然而这一次，我们要以略微不同的方式选取反不变被加项的\trHomotopy{}类，即使得 $\lambda^-(s)=e^{j\sigma(s)}\RR$，其中 $\sigma(0)=\sigma_0$、$\sigma(1)=\sigma_1+k\pi$，而某个 $k\in\ZZ$ 满足 $(-1)^k=w_0w_1$。在 $F_\lambda$ 上选一个\trBraneStructure{}，由一个\trRealValuedPhase{}$\alpha_\lambda^{\#}$ 与一个\trPinStructure{}$P_\lambda^{\#}$ 组成，它们在端点处限制为我们给定的那两个\trLinear\trLagrange\trBrane{}。为实现 \eqref{eq:epsilon-chain-isomorphism}，我们还想要一个\trIsomorphism
\begin{equation*}
f_\lambda\colon P_\lambda^{\#}\longrightarrow P_\lambda^{\#}\otimes\lambda^{\mathrm{top}}(F_\lambda)
\end{equation*}
\trCovering{}$D\iota_M$ 在 $F_\lambda$ 上的作用。\cref{thm:codim-one-iso-coorientations} 告诉我们，这样的\trIsomorphism{}与 $\lambda^-$ 的\trOrientation{}一一对应。我们当然可以调整 $f_\lambda$，使得它在 $s=0$ 处与 \eqref{eq:linear-brane-tensor-iso} 相符；而顺着诱导的\trOrientation{}并利用我们对 $\lambda^-$ 的\trHomotopy{}类所作的假设，即可得知在 $s=1$ 处也是如此。\footnote{译注：原书此处印作“the same thing with hold at $s=1$”；疑应为“will hold”（显见排印笔误）；此处径改。——译者注}作为推论，读出 $o(y)$ 上诱导\trAutomorphism{}的办法是：取 \eqref{eq:absolute-index-orientation-space} 并考察 $D\iota_M$ 在 \eqref{eq:halfplane-index-formula} 上的作用。在我们的情形，这对每一个非零交点 $\lambda^-(s)\cap\lambda^-(1)$（$s<1$）都给出一个 $(-1)$ 因子。这种点的个数将 $\equiv w_0+w_1\pmod{2}$，这就给出了所要的断言（\trEquivalence{}地，这个论证也可以像\cref{thm:equivariant-index-orientation}的证明中那样翻译回\trIndex{}理论的术语）。
\end{proof}
